\subsection{An explicit family robust to all variable orders}
Order robustness is possible in restricted representations without yielding
general circuit lower bounds. Graph-CNF arguments of this kind are established
in knowledge compilation; Bova and Slivovsky
\cite{BovaSlivovsky2017}, Section~4, give substantially stronger
exponential-size examples using expanders. We give an elementary grid example
with a self-contained bound. No novelty is claimed.

Let \(G=(V,E)\) be a fixed finite simple undirected graph with \(m=|V|\).
For a selector string \(y\in\bits^V\), define
\[
 g_G(y)=\bigwedge_{\{u,v\}\in E}\neg(y_u\land y_v).
\]
This is also the satisfiability function, in the selector input \(y\), of
\[
 \Psi_{G,y}(z)=
 \left(\bigwedge_{\{u,v\}\in E}(\neg z_u\lor\neg z_v)\right)
 \land\left(\bigwedge_{v:y_v=1}z_v\right).
\]
Indeed, adjacent selected vertices force an unsatisfied negative clause.
If there are no such vertices, the witness \(z=y\) satisfies every clause.
Thus \(\Psi_{G,y}\) is satisfiable exactly when \(g_G(y)=1\).
The graph and its negative clauses are fixed for this input family; the
\(m\) input bits encode the presence of positive unit clauses.

We retain the deterministic layered model of Lemma~\ref{lem:residual}:
all input bits are read once, in a fixed arbitrary order, with unrestricted
merging of equivalent computations within layers. The measured resource is
the maximum number of reachable states in a layer.

\begin{lemma}[Crossing induced matchings \claimid{L03}]
\label{lem:crossing-matching}
At a cut with read set \(S\subseteq V\), suppose that \(q\) edges
\(\{a_i,b_i\}\), with \(a_i\in S\) and \(b_i\notin S\), form an
induced matching in the \emph{whole graph} \(G\). Every correct layered
diagram for \(g_G\) has at least \(2^q\) reachable states at that cut.
If \(G\) has maximum degree at most \(\Delta\geq1\) and the cut has
\(c\) crossing edges, its width is at least
\(2^{c/(2\Delta^2)}\).
\end{lemma}
\begin{proof}
For each \(A\subseteq\{1,\ldots,q\}\), assign 1 to the read vertices
\(a_i\) with \(i\in A\), and 0 to all other read vertices.
These selected read vertices are independent because the matching is induced
in \(G\). For distinct \(A,A'\), choose \(j\) in their symmetric
difference and set the unread suffix to 1 only at \(b_j\).
There are no edges between endpoints of different matching edges, so the
resulting answer is 0 exactly when \(j\in A\), and likewise for \(A'\).
The \(2^q\) read assignments therefore induce distinct residual functions.
Lemma~\ref{lem:residual} gives the first assertion.

For the second assertion, greedily select a remaining crossing edge
\(\{u,v\}\) and delete every remaining crossing edge incident to
\(N_G(u)\cup N_G(v)\), where these are open neighborhoods in the whole
graph. Both \(u\) and \(v\) belong to that union because \(uv\) is an
edge. The union contains at most \(2\Delta\) vertices, so at most
\(2\Delta^2\) crossing edges are deleted at each selection.
No endpoint of a later selected edge is adjacent to an earlier endpoint;
the selected edges form an induced matching in \(G\), including the absence
of unwanted edges on either side of the cut. When all \(c\) crossing edges
have been deleted, the number selected is at least \(c/(2\Delta^2)\).
Apply the first assertion.
\end{proof}

\begin{proposition}[All-order grid lower bound \claimid{P05}]
\label{prop:grid-orders}
Let \(r\geq2\) be even and let \(G_r\) be the \(r\)-by-\(r\) grid
with horizontal and vertical edges between neighboring vertices.
For every ordering of its \(m=r^2\) selector bits, every correct layered
diagram for \(g_{G_r}\) has width at least \(2^{r/64}\) after the first
\(m/2\) bits. In particular, its maximum width is
\(2^{\Omega(\sqrt m)}\) for every ordering.
\end{proposition}
\begin{proof}
Let \(S\) be any set of \(r^2/2\) grid vertices. Call a row mixed if
it contains vertices both in \(S\) and outside \(S\).
If at least \(r/2\) rows are mixed, each contains a horizontal crossing
edge, giving at least \(r/2\) crossing edges in total.
Otherwise, fewer than \(r/2\) rows are mixed. There must be a row wholly
in \(S\): without one, all \(r^2/2\) vertices of \(S\) would lie in
fewer than \(r/2\) rows, an impossibility. The same argument supplies a
row wholly outside \(S\). In each of the \(r\) columns, a vertical
path between these two rows has a crossing edge. These edges are distinct
across columns, giving at least \(r\) crossing edges.

Thus every such cut has at least \(r/2\) crossing edges. The grid has
maximum degree at most 4. Lemma~\ref{lem:crossing-matching} gives a
crossing induced matching of size at least
\((r/2)/(2\cdot4^2)=r/64\), and hence width at least \(2^{r/64}\).
Every variable order has a prefix of exactly \(r^2/2\) vertices, so the
bound applies to every order.
\end{proof}

This is an explicit family of incidence-encoded SAT instances with an
order-robust lower bound in a specified diagram model. It is not a lower
bound for arbitrary descriptions of SAT tables. The selector-input length
\(m=r^2\) is distinct from the length \(n\) of a conventional binary
formula encoding; the full SAT table in the earlier sections still has
length \(N=2^n\). No bitwise restriction or width-preserving translation
from this incidence encoding to that formula encoding is established here.
The robustness is over all orders of the specified selector bits, not over
all encodings or all computational models.

Moreover, \(g_G\) has a general circuit of size \(O(m+|E|)\), counting
input nodes: test each edge and conjoin the answers. Since
\(|E(G_r)|=2r(r-1)\), this size is \(O(m)\) for the grid family.
The example therefore separates restricted ordered representations from
small general circuits, rather than excluding those circuits.

\begin{table}[htbp]
\centering
% Generated by experiments/graph_order_widths.py; do not edit.
\begin{tabular}{lrr}
\toprule
Family & Input bits & Minimum maximum width \\
\midrule
matching-3 & 6 & 3 \\
cycle-6 & 6 & 5 \\
grid-2x3 & 6 & 4 \\
cube-3 & 8 & 6 \\
grid-3x3 & 9 & 6 \\
\bottomrule
\end{tabular}

\caption{Finite graph-family calculations \claimid{E02}. An exact dynamic
program over subsets computes the minimum, over all variable orders, of
the maximum residual width; it does not enumerate all orders factorially.
These are finite consistency observations, not proofs of the asymptotic bound.}
\label{tab:graph-orders}
\end{table}
